\section{QR 与最小二乘的稳定解法}
\label{sec:08-Numerical-Analysis/05-Numerical-Linear-Algebra/02}

你手上有三个数据点和一个线性模型 \(y=ax+b\)：
\[
(0,1),\quad (1,3),\quad (2,5).
\]
三个点几乎共线，但不可能让一条直线恰好穿过全部三个。一写出来就要面对：
\[
\begin{pmatrix}0&1\\1&1\\2&1\end{pmatrix}
\begin{pmatrix}a\\b\end{pmatrix}
=
\begin{pmatrix}1\\3\\5\end{pmatrix}.
\]
三个方程，两个未知数——方程比未知数多，线性代数课上说这种系统通常没有精确解。那我们应该求什么？

答案是最小二乘：找一个 \(x\) 使得残差向量 \(Ax-b\) 的 2-范数尽可能小。也就是说，在矩阵列空间 \(\mathcal C(A)\) 中找一个离 \(b\) 最近的点 \(A\hat x\)。几何上这件事只有一个意思：残差必须垂直于整个列空间——如果残差在列空间中还有非零投影，沿着那个投影的方向微调一下 \(x\) 就能让残差变得更短，和\"已经最小\"矛盾。

\subsection{正规方程：一个漂亮的公式，一个危险的平方}

把\"残差垂直于列空间\"写成数学条件，就是残差 \(r=b-A\hat x\) 与 \(A\) 的每一列正交：
\[
A^{\trans}r=0.
\]
代入 \(r\) 并移项：
\[
A^{\trans}A\hat x=A^{\trans}b.
\]
这就是正规方程。从推导到最后一步，每一个等号在精确算术中都是完美的。只有两个问题：\(A^{\trans}A\) 的条件数，以及形成 \(A^{\trans}A\) 时发生的相消。

在 2-范数意义下，若 \(A\) 满列秩，
\[
\kappa_2(A^{\trans}A)=\kappa_2(A)^2.
\]
为什么平方？因为 \(A^{\trans}A\) 的特征值是 \(A\) 奇异值的平方——这是线性代数中反复出现的一个事实。如果 \(\kappa_2(A)=10^6\)，那么 \(\kappa_2(A^{\trans}A)=10^{12}\)。原本在双精度下还有大约 10 个有效十进制位的矩阵，经过一次乘法就只剩大约 3 个。而且矩阵乘法中，近共线的列互相减消会使这些丢失直接反映在 \(A^{\trans}A\) 的数值上，而不再只是潜在的理论风险。

正规方程的运算量确实低——但一个小条件数、低精度要求的场合才能放心用。不应该因为公式短就把它当作默认的稳定解法。

\subsection{QR：把几何结构直接变成算法}

回到几何图景。我们想把 \(b\) 投影到 \(\mathcal C(A)\) 上，但 \(A\) 的列不是正交的——直接在这种斜坐标系里做投影很麻烦。如果先把 \(A\) 的列\"摆正\"，投影就退化成简单的坐标提取。

QR 分解做的正是这件事。它把 \(A\) 拆成正交因子和三角因子：
\[
A=QR,
\]
\(Q\) 的列互相正交（若是方阵则正交），\(R\) 是上三角。代入最小二乘目标：
\[
\|Ax-b\|_2
=\|QRx-b\|_2
=\|Rx-Q^{\trans}b\|_2.
\]
第一个等号用了定义，第二个等号是关键：乘一个正交矩阵不改变 2-范数。把 \(Q\) 补成完整的正交方阵后，整个问题拆成两块：前 \(n\) 个坐标必须精确满足
\[
R\hat x=Q^{\trans}b,
\]
而剩余的 \(m-n\) 个坐标就是无论如何也消不掉的残差——它们的大小正好等于最小残差范数。

QR 从头到尾没有形成 \(A^{\trans}A\)，所以不平方条件数。三角系统 \(R\hat x = Q^{\trans}b\) 又可以用前一节说的回代稳定求解。几何上的\"正交投影\"被直接翻译成了数值上稳定的两个操作：正交变换和三角回代。

\subsection{怎样算出 Q 和 R}

有了一个好用的分解，还得有一个不发散的计算方法才能兑现它的理论优势。

经典 Gram--Schmidt 按列处理：把每一列投影到前面已经正交化的方向上，减掉投影，得到新的正交方向。这个过程在列接近平行时会出现严重的相消——减掉一个几乎相等的量后，有用的有效数字几乎全丢了。

Modified Gram--Schmidt 换了一种更新顺序：处理完第 \(k\) 列后，立刻从所有后续列中减掉第 \(k\) 方向上的投影，而不是等到处理那一列时再一次性减去所有前面的投影。这个看似微小的顺序调整显著改善了正交性的保持，但在极端情况下仍可能需要重正交化。

Householder 反射是目前稠密 QR 的标准方法。它不直接构造正交列，而是用一系列正交的反射变换，把矩阵逐列\"按下去\"——每一列处理完后，对角线以下的部分变成零。整个过程等价于左乘一个正交矩阵，而这个正交矩阵不需要显式存储：每个反射只需要一个向量就能完整描述。不显式构造 \(Q\) 既省存储也省正交性被逐渐侵蚀的风险。

Givens 旋转一次只消灭一个元素，用一次二维平面旋转。在稀疏矩阵、流式数据或具有局部结构的矩阵中，消一个元素比搞一整列的反射更灵活。

无论用哪种算法，正交变换的 2-范数保持性质保证了一个关键事实：舍入误差不会因变换本身而被几何放大。这和消元时换行控制增长因子的逻辑是同一回事。

\subsection{近秩亏会给所有方法出难题}

假设设计矩阵的两列分别是 \(t\) 和 \(t+10^{-8}z\)，\(z\) 是和 \(t\) 不相关的随机方向。这两列各自都包含大量信息（模长都很大），但它们几乎完全一样。模型为了拟合两列之间那 \(10^{-8}\) 量级的差异，需要配出两个符号相反、绝对值巨大的系数。训练残差可能很小——模型确实\"解释\"了数据——但系数对数据的最后几位小数点极度敏感。你今天用的数据和明天重新采集的数据，最后几位不一样是很正常的。

QR 避免了正规方程再平方一次条件数的额外惩罚，但它不能把\"两个几乎一样的列\"变成\"两个能分开估计的量\"。这是一个统计意义上的病态，不是数值方法的bug。在面对近共线数据时，应该考虑重新参数化（合并或变换特征）、正则化（用先验信念约束系数大小），或者用 SVD 看清到底是哪些方向被数据几乎确定、哪些方向几乎自由——而不是期待换一个分解就突然给出可信的唯一系数。

列主元 QR
\[
AP=QR
\]
在分解过程中让较独立的列优先进入 \(R\)，可以用于判断\"哪几列组合在一起基本上决定了整个列空间\"——这就是数值秩揭示。SVD 从奇异值的尺度直接给出更稳健的低秩判断，\hyperref[sec:08-Numerical-Analysis/05-Numerical-Linear-Algebra/06]{``SVD、低秩与数值秩''}中会展开。

数值秩不是一个脱离测量单位的绝对整数。判断某个奇异值是否\"等于零\"，需要相对于最大奇异值、矩阵规模、机器精度以及数据噪声大小来设定阈值。把 0.001 在 \(1000\times1000\) 的全 1 矩阵里叫零，和在 \(2\times2\) 的随机矩阵里叫零，含义完全不同。

\subsection{拿到结果后检查什么}

光看残差大小不够，就像上一节里光看残差范数而不看尺度一样不够。一份负责任的最小二乘检查清单至少应包括：

\begin{itemize}
\tightlist
\item
  \(A^{\trans}r\) 是否接近零——这是残差垂直于列空间这一几何条件的直接数值验证；
\item
  \(Q\) 的正交性保持得如何——\(Q^{\trans}Q\) 离单位阵有多远；
\item
  \(R\) 对角元或奇异值是否揭示了近秩亏；
\item
  系数对微小扰动和数据尺度变化是否稳定——给 \(b\) 加个 \(10^{-8}\) 量级的噪声，系数应该变化不大；
\item
  你要求的是秩满时的最小残差解，还是秩亏时的最小范数解——这两种问题有不同答案，QR 和 SVD 的默认行为可能不一样。
\end{itemize}

最小二乘的全部难点可以这样概括：数据本身未能唯一确定答案时，算法可以选择给出一个解，但不能替你决定这个解是否可信。QR 比正规方程更忠实地保留了数据告诉你的那部分信息；至于数据没有告诉你的那部分，任何算法都只是在用某种默认规则填补。
